\makeatletter
\def\input@path{{./}{../../../preamble/}{book2026-2027/preamble/}}
\makeatother
\input{preamble_exams.tex}
\input{graphs/highlights.tex}
% \GTenHighlightsfalse % Uncomment to disable instructional highlights.
% Local semantic roles: indicator, success probability, complementary probability,
% and monetary consequence. Shared rendering and toggle remain authoritative.
\newcommand{\Indicator}[1]{\GTenHLValue{#1}}
\newcommand{\SuccessProb}[1]{\GTenHLProbability{#1}}
\newcommand{\FailureProb}[1]{\GTenHLParameter{#1}}
\newcommand{\Money}[1]{\GTenHLSecond{#1}}
\usepackage{needspace}
\hypersetup{pdftitle={Grade 10 Term 2 Discrete Probability}}
\setlength{\emergencystretch}{2em}
\title{Grade 10, Term 2, Lesson 1, Criterion C3\\Bernoulli Indicators and Financial Consequences}
\author{}
\date{}
\begin{document}
\maketitle
\setcounter{tocdepth}{1}
\setcounter{secnumdepth}{2}
\phantomsection\label{toc}
\tableofcontents
\section*{Evaluated criterion}
\addcontentsline{toc}{section}{Evaluated criterion}
\textbf{C3:} Models binary economic outcomes with Bernoulli random variables, interpreting success and failure in contexts such as default, gain, loss, or claim occurrence.

\textbf{How to practise.} Attempt each C3 task before opening its linked worked solution.
Use the return arrows for the contents, the problem, and its problem list.
All data are hypothetical teaching models. Keep exact expressions where useful;
round final probabilities to four decimal places. Monetary amounts are in USD.

\section{Introduction}\label{sec:t2-c3-introduction}
\subsection*{Highlight legend}
\phantomsection\label{tab:highlight-legend}
\begin{tabularx}{\linewidth}{@{}>{\raggedright\arraybackslash}p{.24\linewidth}>{\raggedright\arraybackslash}p{.25\linewidth}>{\raggedright\arraybackslash}X@{}}
\toprule
\textbf{Symbol or notation} & \textbf{Concept} & \textbf{Brief meaning} \\
\midrule
$\Indicator{X},\ \Indicator{Z}$ & \Indicator{Indicator} & Unitless event label: 1 or 0. \\
$\SuccessProb{p}=P(X=1)$ & \SuccessProb{Success probability} & Chance of the event labelled 1. \\
$\FailureProb{1-p}=P(X=0)$ & \FailureProb{Failure probability} & Chance of the complementary event. \\
$\Money{L},\ \Money{Y}$ & \Money{Monetary}\par\Money{consequence} & Amount of loss, payment, or balance in USD. \\
\bottomrule
\end{tabularx}

The two probability treatments belong to the same binary model but distinguish
success from its complement. Words, symbols, and units carry the meaning even without color.
See the \hyperref[tab:notation-reference]{detailed notation reference} for definitions.

\subsection{Introductory example 1: one invoice}\label{ex:intro-one}
A school enterprise sends one invoice. Let $\Indicator{X}=\Indicator{1}$ if it is paid
by its due date and $\Indicator{X}=\Indicator{0}$ otherwise. Suppose
$\SuccessProb{p}=\SuccessProb{0.90}$. Then the chance of failure (not paid by the due date)
is $\FailureProb{1-p}=1-\SuccessProb{0.90}=\FailureProb{0.10}$.
These two events are mutually exclusive and exhaustive: exactly one happens.
The \Indicator{indicator} records whether the event happens, not the amount paid.
Reversing the labels gives $\Indicator{Z}=1-\Indicator{X}$: now 1 means late payment,
with success probability $\SuccessProb{0.10}$ for this new event.

\subsection{Introductory example 2: one insured event}\label{ex:intro-two}
An insured event occurs during one month with probability $\SuccessProb{0.08}$.
Define $\Indicator{X}=\Indicator{1}$ for occurrence and $\Indicator{X}=\Indicator{0}$ for no occurrence.
Here mathematical \emph{success} is financially undesirable. If occurrence causes a fixed
loss of 200 USD, define $\Money{L}=200\Indicator{X}$ USD. When $X=1$, the
\Money{monetary loss} is $\Money{200}$ USD; when $X=0$, it is $\Money{0}$ USD.
Probabilities and indicators have no monetary units. The event label and its monetary
consequence answer different questions.

\subsection{Binary events, support, and quick checks}
A Bernoulli variable describes one binary observation in a specified period:
success means the event labelled 1; failure means its complement, labelled 0.
Write $X\sim\operatorname{Bernoulli}(p)$, where $\sim$ means ``has the distribution.''
For $0<p<1$, the support (values with positive probability) is $\{0,1\}$.
If $p=0$, only 0 has positive probability; if $p=1$, only 1 does.
At these endpoints use $P(X=1)=p$ and $P(X=0)=1-p$ directly, without evaluating $0^0$.
One observation requires no independence assumption. Counting occurrences across several
observations is a different variable; an indicator of at least one claim does not count claims.

\begin{itemize}
\item \textbf{Check:} Must success be financially good? \emph{Suggested answer:} No; default can be labelled success.
\item \textbf{Check:} If $p=0.08$, what is $P(X=0)$? \emph{Suggested answer:} $1-0.08=0.92$.
\item \textbf{Check:} If $X=1$ means default, what does $X=0$ mean? \emph{Suggested answer:} No default during the stated period.
\item \textbf{Check:} Is the indicator of a 200 USD loss equal to 200? \emph{Suggested answer:} No; it is 1, while the loss is 200 USD.
\end{itemize}
\secInicio

\clearpage
\section{Notation Reference and Practice Procedure}\label{sec:notation-and-procedure}
\subsection{Notation reference}\label{tab:notation-reference}
\begin{tabularx}{\linewidth}{@{}>{\raggedright\arraybackslash}p{.22\linewidth}>{\raggedright\arraybackslash}p{.19\linewidth}>{\raggedright\arraybackslash}p{.23\linewidth}>{\raggedright\arraybackslash}X@{}}
\toprule
\textbf{Symbol or notation} & \textbf{Concept} & \textbf{Definition} & \textbf{Meaning in the introductory examples} \\
\midrule
$\Indicator{X};\ \Indicator{0},\Indicator{1}$ & \Indicator{Indicator} & Unitless random variable; 1 for the defined event, 0 otherwise. &
\textbullet\ \hyperref[ex:intro-one]{Example 1}: on-time payment.\par
\textbullet\ \hyperref[ex:intro-two]{Example 2}: insured-event occurrence. \\
$\SuccessProb{p}=P(X=1)$ & \SuccessProb{Success probability} & Probability of the event labelled 1; $0\leq p\leq1$, no units. &
\textbullet\ \hyperref[ex:intro-one]{Example 1}: $\SuccessProb{0.90}$ for on-time payment.\par
\textbullet\ \hyperref[ex:intro-two]{Example 2}: $\SuccessProb{0.08}$ for occurrence. \\
$\FailureProb{1-p}=P(X=0)$ & \FailureProb{Failure probability} & Probability of the complementary event; the two masses sum to 1. &
\textbullet\ \hyperref[ex:intro-one]{Example 1}: $\FailureProb{0.10}$ for late payment.\par
\textbullet\ \hyperref[ex:intro-two]{Example 2}: $\FailureProb{0.92}$ for no occurrence. \\
$\Indicator{Z}=1-\Indicator{X}$ & \Indicator{Reversed indicator} & Exchanges 0 and 1; $P(Z=1)=1-p$. &
\textbullet\ \hyperref[ex:intro-one]{Example 1}: 1 now means late payment. \\
$\Money{Y}=a\Indicator{X}+b$; $\Money{L}$ & \Money{Monetary}\par\Money{consequence} & Affine transformation; $a,b$ in USD, giving $b$ or $a+b$ USD. &
\textbullet\ \hyperref[ex:intro-two]{Example 2}: $\Money{L}=200\Indicator{X}$ USD; $a=200$, $b=0$. \\
$P(X=x)$; support & Probability mass; possible values & Probability at one value $x$; support lists values with positive probability. &
\textbullet\ \hyperref[ex:intro-one]{Example 1}: $x=1$ means paid on time.\par
\textbullet\ \hyperref[ex:intro-two]{Example 2}: $L$ has support $\{0,200\}$ USD. \\
\bottomrule
\end{tabularx}

\subsection{Reusable practice procedure}
\begin{enumerate}
\item \textbf{Model:} Specify one observation and its period. Define the event labelled 1 and its complement labelled 0. Check that the two labels are exclusive and exhaustive.
\item \textbf{Separate quantities:} Give the unitless indicator and its parameter $p$. Define any monetary variable separately, with USD units. Do not confuse an indicator with a count.
\item \textbf{Calculate:} State $P(X=1)=p$ and $P(X=0)=1-p$. For $0<p<1$, the equivalent probability mass function (PMF) is
\[P(X=x)=p^x(1-p)^{1-x},\qquad x\in\{0,1\}.\]
At $p=0$ or $p=1$, use the direct masses. For a reversed indicator use $Z=1-X$; for money evaluate the transformation at both labels.
\item \textbf{Interpret:} Translate probabilities back into the defined event and period. Explain whether success is financially desirable and distinguish an event probability from an amount of money.
\item \textbf{Check:} Probabilities lie in $[0,1]$ and sum to 1. Check units, supports, and signs. If labels or amounts change, identify exactly which events and probabilities change.
\end{enumerate}
\secInicio[][sec:t2-c3-introduction]

\clearpage
\section{General Worked Problems}\label{sec:t2-c3-general}
\phantomsection\label{sec:t2-c3-foundational}\label{sec:t2-c3-intermediate}
Each task addresses C3 by defining and interpreting a binary economic outcome.
The general problems establish the model; the foundational and intermediate problems apply it.
\SectionProblemsTableOfContents{
    \SectionProblemLink{prob:c3-binary}{Introductory general problem: one binary trial} \\
    \SectionProblemLink{prob:c3-reverse}{Advanced general problem: reversing success and measuring money} \\
    \SectionProblemLink{prob:c3-loan}{Foundational problem: one loan default} \\
    \SectionProblemLink{prob:c3-claim}{Foundational problem: one policy claim} \\
    \SectionProblemLink{prob:c3-invoice}{Intermediate problem: invoice payment and reversal} \\
    \SectionProblemLink{prob:c3-gain}{Intermediate problem: investment gain}
}

\subsection{Introductory general problem: one binary trial}\label{prob:c3-binary}
A school enterprise observes one invoice at its due date. Let A be the event that it is paid on time, with probability 0.3. Define $X=1$ if A occurs and 0 otherwise.

\QuestionsTasks[sec:t2-c3-general]{
    {3}/{sol:c3-binary}/{Give the support and calculate both masses from the Bernoulli formula. If $p$ changes to 0 or 1, describe the masses and support without using $0^0$.}
}

\subsubsection*{Complete worked solution}\label{sol:c3-binary}
\par\noindent\begin{minipage}{\linewidth}
\textbf{Step 1: Define and model.} $\Indicator{X}$ indicates A; support $\{0,1\}$, no monetary units, parameter $\SuccessProb{p}=\SuccessProb{0.3}$. A and its complement are exclusive and exhaustive, so the Bernoulli model applies.
\end{minipage}\par\medskip
\par\noindent\begin{minipage}{\linewidth}
\textbf{Step 2: Write the formula.} $P(X=x)=p^x(1-p)^{1-x}$ for $x\in\{0,1\}$. The notation $x\in\{0,1\}$ means x is either 0 or 1; a nonzero number to power zero is 1.
\end{minipage}\par\medskip
\par\noindent\begin{minipage}{\linewidth}
\textbf{Step 3: Substitute.} $P(\Indicator{X}=\Indicator{1})=\SuccessProb{0.3}^1(\FailureProb{0.7})^0=\SuccessProb{0.3000}$, and $P(\Indicator{X}=\Indicator{0})=\SuccessProb{0.3}^0(\FailureProb{0.7})^1=\FailureProb{0.7000}$.
\end{minipage}\par\medskip
\par\noindent\begin{minipage}{\linewidth}
\textbf{Step 4: Check and interpret.} $0.3+0.7=1$. The value 1 labels on-time payment, not the invoice amount; its probability is 0.3000.
\end{minipage}\par\medskip

\textbf{Follow-up reasoning.} For the same event A, the direct rules $P(\Indicator{X}=\Indicator{1})=p$, $P(\Indicator{X}=\Indicator{0})=1-p$ apply at the endpoints. If $p=0$, the masses are 0 and 1 and the support is $\{0\}$. If $p=1$, they are 1 and 0 and the support is $\{1\}$. In each case the masses total 1: occurrence is impossible or certain, respectively.
\secInicio[prob:c3-binary][sec:t2-c3-general]

\clearpage
\subsection{Advanced general problem: reversing success and measuring money}\label{prob:c3-reverse}
For one financial contract at settlement, $X$ indicates an adverse event of probability 0.2: $X=1$ for occurrence and $X=0$ otherwise. The balance is 12 USD before a 50 USD deduction if the event occurs. Define $Z=1-X$ and $Y=12-50X$ USD.

\QuestionsTasks[sec:t2-c3-general]{
    {3}/{sol:c3-reverse}/{Find both indicator PMFs, the monetary support and $P(Y<0)$. Explain whether $Y$ is Bernoulli. Compare with an adverse-event probability of 0.35, leaving the monetary amounts fixed.}
}

\subsubsection*{Complete worked solution}\label{sol:c3-reverse}
\par\noindent\begin{minipage}{\linewidth}
\textbf{Step 1: Define.} $\Indicator{X}=\Indicator{1}$ labels the adverse event; $\Indicator{Z}=\Indicator{1}$ labels its absence. Both have support $\{0,1\}$ and group the same two exclusive, exhaustive outcomes, so both are Bernoulli. $\Money{Y}$ measures money, with support $\{-38,12\}$ USD.
\end{minipage}\par\medskip
\par\noindent\begin{minipage}{\linewidth}
\textbf{Step 2: Model and reverse.} $P(\Indicator{X}=\Indicator{1})=p=0.2$ and $P(\Indicator{X}=\Indicator{0})=1-p=0.8$, so $X\sim\operatorname{Bernoulli}(0.2)$, $Z\sim\operatorname{Bernoulli}(0.8)$. One event is being observed, so no trial independence assumption is needed.
\end{minipage}\par\medskip
\par\noindent\begin{minipage}{\linewidth}
\textbf{Step 3: Calculate.} $P(\Indicator{Z}=\Indicator{1})=1-p=1-0.2=0.8$, $P(\Indicator{Z}=\Indicator{0})=0.2$. Also $Y(1)=12-50=-38$ and $Y(0)=12$. Hence $P(Y<0)=P(\Indicator{X}=\Indicator{1})=\SuccessProb{0.2000}$; $P(\Money{Y}=\Money{12})=0.8$.
\end{minipage}\par\medskip
\par\noindent\begin{minipage}{\linewidth}
\textbf{Step 4: Check and interpret.} Each PMF totals 1. $\Money{Y}$ is a monetary transformation of an indicator, but it is not itself Bernoulli because its values are not 0 and 1. Reversing labels does not change the event causing the loss.
\end{minipage}\par\medskip

\textbf{Follow-up reasoning.} The event and transformation stay the same. With $p=0.35$, $P(\Indicator{X}=\Indicator{1})=\SuccessProb{0.35}$, $P(\Indicator{X}=\Indicator{0})=\FailureProb{0.65}$, and $P(\Indicator{Z}=\Indicator{1})=1-0.35=0.65$. The monetary values remain $-38$ and 12 USD, but $P(Y<0)=0.3500$. Raising the adverse-event probability raises the chance of a loss; reversing labels alone never changes that chance.
\secInicio[prob:c3-reverse][sec:t2-c3-general]

\clearpage
\subsection{Foundational problem: one loan default}\label{prob:c3-loan}
One loan has a hypothetical probability 0.04 of default within one year. Default causes a fixed loss of 800 USD; otherwise default loss is zero.

\QuestionsTasks[sec:t2-c3-general]{
    {3}/{sol:c3-loan}/{Define the indicator of default within one year; give its PMF and the financial interpretation. If the fixed default loss rises to 1000 USD with the same default probability, explain what changes.}
}

\subsubsection*{Complete worked solution}\label{sol:c3-loan}
\par\noindent\begin{minipage}{\linewidth}
\textbf{Step 1: Define the event.} $\Indicator{X}=\Indicator{1}$ if the loan defaults within one year, and $\Indicator{X}=\Indicator{0}$ otherwise. Support is $\{0,1\}$, without monetary units; $\SuccessProb{p}=\SuccessProb{0.04}$.
\end{minipage}\par\medskip
\par\noindent\begin{minipage}{\linewidth}
\textbf{Step 2: Justify the model.} This is one observation with two exclusive, exhaustive labels, so it is Bernoulli. Repeated-observation independence is not needed here.
\end{minipage}\par\medskip
\par\noindent\begin{minipage}{\linewidth}
\textbf{Step 3: Translate and calculate.} $P(X=x)=p^x(1-p)^{1-x}$ for $x=0,1$. Thus $P(\Indicator{X}=\Indicator{1})=\SuccessProb{0.04}$, $P(\Indicator{X}=\Indicator{0})=1-\SuccessProb{0.04}=\FailureProb{0.96}$. $\Money{L}=\Money{800}\Indicator{X}$ USD: $L(1)=800(1)=\Money{800}$ and $L(0)=800(0)=\Money{0}$. Thus $P(\Money{L}=\Money{800})=P(\Indicator{X}=\Indicator{1})=\SuccessProb{0.04}$ and $P(\Money{L}=\Money{0})=0.96$. Mathematical success is financially undesirable.
\end{minipage}\par\medskip
\par\noindent\begin{minipage}{\linewidth}
\textbf{Step 4: Check and interpret.} The non-negative masses sum to $ 0.04+0.96=1$. The \SuccessProb{success probability} is $\SuccessProb{0.0400}$. The \Indicator{indicator} records occurrence; the \Money{monetary variable} measures its consequence in USD.
\end{minipage}\par\medskip

\textbf{Follow-up reasoning.} The event remains default within one year, so $X\sim\operatorname{Bernoulli}(0.04)$ is unchanged. The new loss is $\Money{L}=\Money{1000}X$ USD: $L(1)=1000$ and $L(0)=0$. Thus $P(\Money{L}=\Money{1000})=0.04$ and $P(\Money{L}=\Money{0})=0.96$. The loss size rises; its probability remains 0.0400.
\secInicio[prob:c3-loan][sec:t2-c3-general]

\clearpage
\subsection{Foundational problem: one policy claim}\label{prob:c3-claim}
One insurance policy has probability 0.12 of at least one claim this month. The simplified contract pays 300 USD once if any claim occurs.

\QuestionsTasks[sec:t2-c3-general]{
    {3}/{sol:c3-claim}/{Define the indicator of at least one claim this month; give its PMF and the financial interpretation. Could this indicator tell whether there were one or three claims? Explain whether $Y=300X$ still works if the contract pays 300 USD per claim instead.}
}

\subsubsection*{Complete worked solution}\label{sol:c3-claim}
\par\noindent\begin{minipage}{\linewidth}
\textbf{Step 1: Define the event.} $\Indicator{X}=\Indicator{1}$ if at least one claim occurs this month, and $\Indicator{X}=\Indicator{0}$ otherwise. Support is $\{0,1\}$, without monetary units; $\SuccessProb{p}=\SuccessProb{0.12}$.
\end{minipage}\par\medskip
\par\noindent\begin{minipage}{\linewidth}
\textbf{Step 2: Justify the model.} This is one observation with two exclusive, exhaustive labels, so it is Bernoulli. Repeated-observation independence is not needed here.
\end{minipage}\par\medskip
\par\noindent\begin{minipage}{\linewidth}
\textbf{Step 3: Translate and calculate.} $P(X=x)=p^x(1-p)^{1-x}$ for $x=0,1$. Thus $P(\Indicator{X}=\Indicator{1})=\SuccessProb{0.12}$, $P(\Indicator{X}=\Indicator{0})=1-\SuccessProb{0.12}=\FailureProb{0.88}$. $\Money{Y}=\Money{300}\Indicator{X}$ USD: $Y(1)=300(1)=\Money{300}$ and $Y(0)=300(0)=\Money{0}$. Hence $P(Y>0)=P(\Indicator{X}=\Indicator{1})=\SuccessProb{0.12}$. The indicator does not count claims.
\end{minipage}\par\medskip
\par\noindent\begin{minipage}{\linewidth}
\textbf{Step 4: Check and interpret.} The non-negative masses sum to $ 0.12+0.88=1$. The \SuccessProb{success probability} is $\SuccessProb{0.1200}$. The \Indicator{indicator} records occurrence; the \Money{monetary variable} measures its consequence in USD.
\end{minipage}\par\medskip

\textbf{Follow-up reasoning.} Both one and three claims give $\Indicator{X}=\Indicator{1}$, so this binary record cannot distinguish their counts. Under the original once-only payment rule, $\Money{Y}=\Money{300}X$ is valid. Under a per-claim rule, let $N$ be the number of claims: the payment would be $300N$ USD, not generally $300X$. The given $P(N\geq1)=0.12$ determines the indicator PMF but does not determine the probabilities of individual claim counts or payment amounts.
\secInicio[prob:c3-claim][sec:t2-c3-general]

\clearpage
\subsection{Intermediate problem: invoice payment and reversal}\label{prob:c3-invoice}
One invoice is paid by its due date with probability 0.85. Late payment triggers a 6 USD fee.

\QuestionsTasks[sec:t2-c3-general]{
    {3}/{sol:c3-invoice}/{Define the indicator of on-time invoice payment; give its PMF and the financial interpretation. Reverse success to lateness and find the fee probability. Compare an on-time probability of 0.95 with the original case, keeping the fee at 6 USD.}
}

\subsubsection*{Complete worked solution}\label{sol:c3-invoice}
\par\noindent\begin{minipage}{\linewidth}
\textbf{Step 1: Define the event.} $\Indicator{X}=\Indicator{1}$ if the invoice is paid on time, and $\Indicator{X}=\Indicator{0}$ otherwise. Support is $\{0,1\}$, without monetary units; $\SuccessProb{p}=\SuccessProb{0.85}$.
\end{minipage}\par\medskip
\par\noindent\begin{minipage}{\linewidth}
\textbf{Step 2: Justify the model.} This is one observation with two exclusive, exhaustive labels, so it is Bernoulli. Repeated-observation independence is not needed here.
\end{minipage}\par\medskip
\par\noindent\begin{minipage}{\linewidth}
\textbf{Step 3: Translate and calculate.} $P(X=x)=p^x(1-p)^{1-x}$ for $x=0,1$. Thus $P(\Indicator{X}=\Indicator{1})=\SuccessProb{0.85}$, $P(\Indicator{X}=\Indicator{0})=1-\SuccessProb{0.85}=\FailureProb{0.15}$. For lateness use $\Indicator{Z}=\Indicator{1}-X$, with parameter $1-p=0.15$. The fee is $\Money{L}=\Money{6}Z$ USD, so $P(\Money{L}=\Money{6})=P(\Indicator{Z}=\Indicator{1})=0.15$.
\end{minipage}\par\medskip
\par\noindent\begin{minipage}{\linewidth}
\textbf{Step 4: Check and interpret.} The non-negative masses sum to $ 0.85+0.15=1$. The \SuccessProb{success probability} is $\SuccessProb{0.8500}$. The \Indicator{indicator} records occurrence; the \Money{monetary variable} measures its consequence in USD.
\end{minipage}\par\medskip

\textbf{Follow-up reasoning.} The event is still on-time payment, with $P(\Indicator{X}=\Indicator{1})=\SuccessProb{0.95}$, $P(\Indicator{X}=\Indicator{0})=1-\SuccessProb{0.95}=\FailureProb{0.05}$. The lateness indicator has $P(\Indicator{Z}=\Indicator{1})=0.05$. Hence $\Money{L}=\Money{6}Z$ USD still has values 0 and 6, but $P(\Money{L}=\Money{6})=0.0500$, down from 0.1500 by 0.10 (10 percentage points). The fee size is unchanged.
\secInicio[prob:c3-invoice][sec:t2-c3-general]

\clearpage
\subsection{Intermediate problem: investment gain}\label{prob:c3-gain}
One investment yields 15 USD with probability 0.6 and loses 10 USD otherwise.

\QuestionsTasks[sec:t2-c3-general]{
    {3}/{sol:c3-gain}/{Define the indicator of a positive investment gain; give its PMF and the financial interpretation. If the positive outcome falls from 15 USD to 5 USD, with the loss and probabilities unchanged, give the new monetary formula and explain the effect on the gain indicator.}
}

\subsubsection*{Complete worked solution}\label{sol:c3-gain}
\par\noindent\begin{minipage}{\linewidth}
\textbf{Step 1: Define the event.} $\Indicator{X}=\Indicator{1}$ if the investment has a positive gain, and $\Indicator{X}=\Indicator{0}$ otherwise. Support is $\{0,1\}$, without monetary units; $\SuccessProb{p}=\SuccessProb{0.6}$.
\end{minipage}\par\medskip
\par\noindent\begin{minipage}{\linewidth}
\textbf{Step 2: Justify the model.} This is one observation with two exclusive, exhaustive labels, so it is Bernoulli. Repeated-observation independence is not needed here.
\end{minipage}\par\medskip
\par\noindent\begin{minipage}{\linewidth}
\textbf{Step 3: Translate and calculate.} $P(X=x)=p^x(1-p)^{1-x}$ for $x=0,1$. Thus $P(\Indicator{X}=\Indicator{1})=\SuccessProb{0.6}$, $P(\Indicator{X}=\Indicator{0})=1-\SuccessProb{0.6}=\FailureProb{0.4}$. $\Money{G}=\Money{-10}+25\Indicator{X}$ USD: $G(1)=-10+25=\Money{15}$ and $G(0)=-10+0=\Money{-10}$. Therefore $P(\Money{G}=\Money{15})=0.6$, $P(\Money{G}=\Money{-10})=0.4$ and $P(G<0)=0.4$. The indicator value 1 does not mean 1 USD.
\end{minipage}\par\medskip
\par\noindent\begin{minipage}{\linewidth}
\textbf{Step 4: Check and interpret.} The non-negative masses sum to $ 0.6+0.4=1$. The \SuccessProb{success probability} is $\SuccessProb{0.6000}$. The \Indicator{indicator} records occurrence; the \Money{monetary variable} measures its consequence in USD.
\end{minipage}\par\medskip

\textbf{Follow-up reasoning.} The event remains positive gain, so $P(\Indicator{X}=\Indicator{1})=\SuccessProb{0.6}$ and $P(\Indicator{X}=\Indicator{0})=\FailureProb{0.4}$ are unchanged. The new monetary formula is $\Money{G}=\Money{-10}+15X$ USD: $G(1)=5$, $G(0)=-10$. Thus the gain size falls, while $P(G>0)=0.6000$ and $P(G<0)=0.4000$ stay the same. Both monetary PMFs total 1.
\secInicio[prob:c3-gain][sec:t2-c3-general]
\clearpage
\section{Advanced Problems}\label{sec:t2-c3-advanced}
These C3 tasks compare labels, monetary consequences, and binary groupings of underlying outcomes.
\SectionProblemsTableOfContents{
    \SectionProblemLink{prob:c3-chargeback}{Chargeback and merchant balance} \\
    \SectionProblemLink{prob:c3-categories}{Three sales categories}
}

\subsection{Chargeback and merchant balance}\label{prob:c3-chargeback}
One sale brings 30 USD initially. With probability 0.08 a chargeback removes 30 USD and adds a 5 USD fee.

\QuestionsTasks[sec:t2-c3-advanced]{
    {3}/{sol:c3-chargeback}/{Define the indicator of a chargeback on the sale; give its PMF and the financial interpretation. Find the negative-balance probability and the reversed parameter. Compare the balance using the reversed indicator. Then remove the 5 USD fee, keeping the chargeback probability unchanged: does the negative-balance probability stay the same?}
}

\subsubsection*{Complete worked solution}\label{sol:c3-chargeback}
\par\noindent\begin{minipage}{\linewidth}
\textbf{Step 1: Define the event.} $\Indicator{X}=\Indicator{1}$ if the sale produces a chargeback, and $\Indicator{X}=\Indicator{0}$ otherwise. Support is $\{0,1\}$, without monetary units; $\SuccessProb{p}=\SuccessProb{0.08}$.
\end{minipage}\par\medskip
\par\noindent\begin{minipage}{\linewidth}
\textbf{Step 2: Justify the model.} This is one observation with two exclusive, exhaustive labels, so it is Bernoulli. Repeated-observation independence is not needed here.
\end{minipage}\par\medskip
\par\noindent\begin{minipage}{\linewidth}
\textbf{Step 3: Translate and calculate.} $P(X=x)=p^x(1-p)^{1-x}$ for $x=0,1$. Thus $P(\Indicator{X}=\Indicator{1})=\SuccessProb{0.08}$, $P(\Indicator{X}=\Indicator{0})=1-\SuccessProb{0.08}=\FailureProb{0.92}$. $\Money{Y}=\Money{30}-35X$ USD, with values 30 and $-5$. Thus $P(Y<0)=P(\Indicator{X}=\Indicator{1})=\SuccessProb{0.08}$. Reversing success to no chargeback gives parameter 0.92, but the loss probability stays 0.08.
\end{minipage}\par\medskip
\par\noindent\begin{minipage}{\linewidth}
\textbf{Step 4: Check and interpret.} The non-negative masses sum to $ 0.08+0.92=1$. The \SuccessProb{success probability} is $\SuccessProb{0.0800}$. The \Indicator{indicator} records occurrence; the \Money{monetary variable} measures its consequence in USD.
\end{minipage}\par\medskip

\textbf{Follow-up reasoning.} Let $\Indicator{Z}=\Indicator{1}-X$ indicate no chargeback. Substituting $\Indicator{X}=\Indicator{1}-Z$ gives $\Money{Y}=\Money{30}-35(1-Z)=-5+35Z$ USD, so $Y<0$ means $\Indicator{Z}=\Indicator{0}$ and has probability 0.08. Without the fee the balance becomes $Y_0=30-30X=30Z$ USD, with values 0 and 30. Thus $P(Y_0<0)=0$, although the chargeback probability stays 0.08. A chargeback now produces zero balance, not a negative balance.
\secInicio[prob:c3-chargeback][sec:t2-c3-advanced]

\clearpage
\subsection{Three sales categories}\label{prob:c3-categories}
One enquiry produces no sale, a basic sale or a premium sale with probabilities $0.5,0.3,0.2$. Basic revenue is 20 USD; premium revenue is 50 USD.

\QuestionsTasks[sec:t2-c3-advanced]{
    {3}/{sol:c3-categories}/{Explain why revenue is not Bernoulli. Give PMFs for purchase and premium indicators, and find the chance of revenue at least 20 USD. Form the no-purchase indicator and express revenue using both purchase and premium indicators. If basic revenue rises to 25 USD with the same category probabilities, explain what changes.}
}

\subsubsection*{Complete worked solution}\label{sol:c3-categories}
\par\noindent\begin{minipage}{\linewidth}
\textbf{Step 1: Define.} $\Money{R}$ is revenue in USD, support $\{0,20,50\}$. Define $\Indicator{X}=\Indicator{1}$ for any purchase and 0 otherwise; $\Indicator{Z}=\Indicator{1}$ for premium and 0 otherwise. Both indicators have support $\{0,1\}$.
\end{minipage}\par\medskip
\par\noindent\begin{minipage}{\linewidth}
\textbf{Step 2: Select.} Three revenue values exclude Bernoulli for $\Money{R}$. Each indicator groups the exclusive, exhaustive outcomes into two categories, so each is Bernoulli. They concern the same enquiry and are not assumed independent.
\end{minipage}\par\medskip
\par\noindent\begin{minipage}{\linewidth}
\textbf{Step 3: Calculate.} For exclusive categories, add their probabilities; for any indicator $I$, $P(I=0)=1-P(I=1)$. $P(\Indicator{X}=\Indicator{1})=P(\text{basic})+P(\text{premium})=0.3+0.2=0.5$; $P(\Indicator{X}=\Indicator{0})=\FailureProb{0.5}$. Also $P(\Indicator{Z}=\Indicator{1})=0.2$, $P(\Indicator{Z}=\Indicator{0})=0.8$. Therefore $P(R\geq20)=P(\Indicator{X}=\Indicator{1})=\SuccessProb{0.5000}$.
\end{minipage}\par\medskip
\par\noindent\begin{minipage}{\linewidth}
\textbf{Step 4: Check and interpret.} Each PMF totals 1. Premium implies purchase, but the two indicators answer different questions. Revenue measures monetary size.
\end{minipage}\par\medskip

\textbf{Follow-up reasoning.} Let $\Indicator{W}=\Indicator{1}-X$ indicate no purchase. Then $P(\Indicator{W}=\Indicator{1})=1-0.5=0.5$ and $P(\Indicator{W}=\Indicator{0})=0.5$. The formula $\Money{R}=\Money{20}X+30Z$ USD gives 0 for no sale, 20 for basic, and 50 for premium; $\Indicator{Z}=\Indicator{1}$ implies $\Indicator{X}=\Indicator{1}$, so the combination $(X,Z)=(0,1)$ is impossible. With basic revenue 25 USD, $\Money{R}=\Money{25}X+25Z$ USD gives 0, 25, and 50. The category and indicator probabilities stay unchanged; $P(R\geq20)=0.5000$ still, but $P(R\geq25)$ rises from 0.2000 to 0.5000. Revenue remains a three-valued monetary variable, not a Bernoulli indicator.
\secInicio[prob:c3-categories][sec:t2-c3-advanced]

\subsection*{Closing reminders: key ideas and common errors}
Define success before assigning $p$; it may mean default or loss. Reversing the label
changes its parameter to $1-p$, not the underlying financial outcome. Distinguish
an indicator, a monetary consequence, and a count across observations.
A monetary variable is not Bernoulli just because it has two possible amounts:
a Bernoulli variable uses the labels 0 and 1. Several underlying categories may still
be grouped into one binary event. A single Bernoulli observation requires no independence
assumption. At endpoint probabilities use direct masses, not $0^0$.
\secInicio[][sec:t2-c3-advanced]
\end{document}
